\documentclass[12pt]{article}

\usepackage{amsmath, amssymb, amsthm}
\usepackage{geometry}
\usepackage{hyperref}
\usepackage{booktabs}
\usepackage{array}

\geometry{margin=1in}

\theoremstyle{definition}
\newtheorem{definition}{Definition}[section]
\newtheorem{example}[definition]{Example}
\newtheorem{remark}[definition]{Remark}

\theoremstyle{plain}
\newtheorem{theorem}[definition]{Theorem}
\newtheorem{proposition}[definition]{Proposition}
\newtheorem{conjecture}[definition]{Conjecture}

% Shortcuts
\newcommand{\RR}{\mathbb{R}}
\newcommand{\CC}{\mathbb{C}}
\newcommand{\HH}{\mathbb{H}}
\newcommand{\OO}{\mathbb{O}}
\newcommand{\Sp}{\mathrm{Sp}}
\newcommand{\SU}{\mathrm{SU}}
\newcommand{\SO}{\mathrm{SO}}
\newcommand{\Spin}{\mathrm{Spin}}
\newcommand{\U}{\mathrm{U}}
\newcommand{\Hol}{\mathrm{Hol}}

\title{K\"{a}hler, Hyperk\"{a}hler, and Quaternionic Geometry\\
\large Notes from the nLab}
\author{}
\date{}

\begin{document}

\maketitle
\tableofcontents

\newpage

%------------------------------------------------------------
\section{Overview: The Normed Division Algebra Hierarchy}
%------------------------------------------------------------

The four normed division algebras $\RR, \CC, \HH, \OO$ organize the landscape of special holonomy Riemannian manifolds. In Berger's classification, each algebra corresponds to a family of Riemannian manifolds whose holonomy group is a subgroup of the relevant structure group:

\begin{center}
\renewcommand{\arraystretch}{1.4}
\begin{tabular}{cccc}
\toprule
\textbf{Algebra} $\mathbb{A}$ & \textbf{Riemannian $\mathbb{A}$-manifold} & \textbf{Special (restricted holonomy)} \\
\midrule
$\RR$ & Riemannian manifold & Oriented Riemannian manifold \\
$\CC$ & K\"{a}hler manifold & Calabi--Yau manifold \\
$\HH$ & Quaternion-K\"{a}hler manifold & Hyperk\"{a}hler manifold \\
$\OO$ & $\Spin(7)$-manifold & $G_2$-manifold \\
\bottomrule
\end{tabular}
\end{center}

\medskip

This hierarchy is reflected in Berger's theorem on the classification of irreducible, non-symmetric Riemannian holonomy groups.

%------------------------------------------------------------
\section{K\"{a}hler Manifolds}
%------------------------------------------------------------

\begin{definition}
A \textbf{K\"{a}hler manifold} is a Riemannian manifold $(M, g)$ of real dimension $2n$ equipped with an integrable complex structure $J$ such that:
\begin{itemize}
  \item $J$ is compatible with $g$: the pairing $\omega(X, Y) := g(JX, Y)$ is skew-symmetric, and
  \item $\omega$ is closed: $d\omega = 0$.
\end{itemize}
Equivalently, the holonomy group $\Hol(M, g)$ is a subgroup of $\U(n)$.
\end{definition}

The 2-form $\omega$ is called the \textbf{K\"{a}hler form}. K\"{a}hler manifolds simultaneously carry Riemannian, symplectic, and complex structures, all mutually compatible.

\begin{definition}
A \textbf{Calabi--Yau manifold} is a K\"{a}hler manifold whose holonomy reduces further to $\SU(n) \subset \U(n)$. Equivalently, it is a K\"{a}hler manifold with a holomorphic volume form (trivial canonical bundle).
\end{definition}

\begin{remark}
By the Hodge decomposition theorem, odd-degree cohomology of a compact K\"{a}hler manifold is always even-dimensional. This provides a useful obstruction: the Hopf surface $S^1 \times S^3$ is diffeomorphic to a product with odd first Betti number, hence cannot be K\"{a}hler.
\end{remark}

%------------------------------------------------------------
\section{Hyperk\"{a}hler Manifolds}
%------------------------------------------------------------

\begin{definition}
A \textbf{hyperk\"{a}hler manifold} is a Riemannian manifold $(M, g)$ of real dimension $4n$ equipped with three integrable complex structures $I, J, K$ each K\"{a}hler with respect to $g$, satisfying the quaternionic identities:
\[
I^2 = J^2 = K^2 = IJK = -1.
\]
Equivalently, the holonomy group $\Hol(M, g)$ is a subgroup of $\Sp(n)$.
\end{definition}

At each point $x \in M$, the tangent space $T_x M$ is a quaternionic vector space isomorphic to $\HH^n$. The group $\Sp(n)$ acts on $\HH^n$ by left-multiplication by quaternionic unitary matrices, preserving $I$, $J$, and $K$.

\begin{proposition}
Every hyperk\"{a}hler manifold is Ricci-flat, hence is a Calabi--Yau manifold (when viewed as a complex manifold via any one of the complex structures $I, J, K$).
\end{proposition}

\begin{proposition}
A hyperk\"{a}hler manifold, considered as a complex manifold via $I$, is holomorphically symplectic: it admits a holomorphic, non-degenerate, closed $2$-form $\Omega = \omega_J + i\,\omega_K$.
\end{proposition}

\subsection{Preserved Differential Forms}

The three K\"{a}hler forms $\omega_1, \omega_2, \omega_3$ associated to $I, J, K$ respectively are each parallel. The sphere
\[
\omega = a\,\omega_1 + b\,\omega_2 + c\,\omega_3, \quad a^2 + b^2 + c^2 = 1,
\]
is a whole $S^2$ worth of K\"{a}hler structures on $(M, g)$.

\subsection{Compact Hyperk\"{a}hler Manifolds}

The classification of compact hyperk\"{a}hler manifolds is largely open. The known examples are:

\begin{example}
The only known compact hyperk\"{a}hler manifolds (up to deformation) are:
\begin{enumerate}
  \item Hilbert schemes of points $X^{[n+1]}$ for $X$ a K3 surface (Beauville 1983),
  \item Generalized Kummer varieties $(T^4)^{[n]} \to T^4$ for $X = T^4$ a 4-torus (Beauville 1983),
  \item Two exceptional examples due to O'Grady in dimensions 10 and 6 (O'Grady 1999, 2003).
\end{enumerate}
\end{example}

\subsection{Rozansky--Witten Weight Systems}

Every hyperk\"{a}hler manifold $M$ induces a Rozansky--Witten weight system with coefficients in certain Dolbeault cohomology groups $H^{0,*}(M, \Lambda^* T^{1,0}M)$. For compact hyperk\"{a}hler manifolds, these descend to actual numerical weight systems valued in the ground field.

\subsection{Physics: Coulomb and Higgs Branches}

Both the Coulomb branch and the Higgs branch of $D = 3$, $\mathcal{N} = 4$ super Yang--Mills theories are hyperk\"{a}hler manifolds (Seiberg--Witten 1996). In special cases they are compact.

%------------------------------------------------------------
\section{Quaternion-K\"{a}hler Manifolds}
%------------------------------------------------------------

\begin{definition}
A Riemannian manifold $(X, g)$ of dimension $4n$ (with $n \geq 2$) is a \textbf{quaternion-K\"{a}hler manifold} if its holonomy group is a subgroup of
\[
\Sp(n) \cdot \Sp(1) := (\Sp(n) \times \Sp(1))\,/\,(\mathbb{Z}/2),
\]
where $\Sp(1) \cong \SU(2) \cong \Spin(3)$ and the $\mathbb{Z}/2$ quotient identifies the diagonal centers.
\end{definition}

\begin{remark}
If the holonomy further reduces to the $\Sp(n)$-factor alone, one recovers a hyperk\"{a}hler manifold. Thus hyperk\"{a}hler manifolds are the ``special'' quaternion-K\"{a}hler manifolds, in parallel with how Calabi--Yau manifolds are special K\"{a}hler manifolds.
\end{remark}

The holonomy group $\Sp(n) \cdot \Sp(1)$ does not globally preserve any particular complex structure on $TX$, but does preserve a rank-3 subbundle of $\mathrm{End}(TX)$ locally spanned by $\{I, J, K\}$ at each point. The quaternionic structure is thus defined only up to $\Sp(1)$-rotation in the fiber.

\subsection{Preserved Differential Form}

A quaternion-K\"{a}hler manifold carries a globally defined parallel $4$-form:
\[
\omega_4 = \omega_1 \wedge \omega_1 + \omega_2 \wedge \omega_2 + \omega_3 \wedge \omega_3,
\]
where $\omega_1, \omega_2, \omega_3$ are the local K\"{a}hler forms associated to $I, J, K$. This $4$-form is well-defined globally despite $I, J, K$ only being locally defined.

\subsection{Einstein Property}

\begin{proposition}
Every quaternion-K\"{a}hler manifold is an Einstein manifold. In particular, the scalar curvature $R$ is constant.
\end{proposition}

This is a consequence of the holonomy reduction and is a strict constraint: the scalar curvature $R \in \RR$ partitions quaternion-K\"{a}hler manifolds into three cases: $R > 0$, $R = 0$, and $R < 0$.

\subsection{Relation to Hyperk\"{a}hler}

\begin{proposition}
A quaternion-K\"{a}hler manifold $(X, g)$ is hyperk\"{a}hler if and only if $R(g) = 0$.
\end{proposition}

Thus hyperk\"{a}hler manifolds are precisely the Ricci-flat (equivalently, scalar-flat) quaternion-K\"{a}hler manifolds.

\subsection{As Quaternionic Manifolds}

\begin{proposition}
Every quaternion-K\"{a}hler manifold is automatically a quaternionic manifold. Specifically, the Levi-Civita connection extends to a torsion-free connection on the $\Sp(n) \cdot \Sp(1)$-holonomy bundle.
\end{proposition}

Such an extension $\nabla_{\mathrm{quat}}$ is not unique: any tensor $\mathcal{S}$ valued in the first prolongation of the structure Lie algebra yields another valid connection $\nabla_{\mathrm{quat}} + \mathcal{S}$.

\subsection{Characteristic Classes (Dimension 8)}

\begin{proposition}
Let $X$ be a closed smooth 8-manifold with a Spin structure and a $G$-structure for $G = \Sp(2) \cdot \Sp(1) \hookrightarrow \SO(8)$. Then the Euler class $\chi$, second Pontryagin class $p_2$, and the cup-square of the first Pontryagin class $(p_1)^2$ satisfy:
\[
8\chi = 4p_2 - (p_1)^2.
\]
\end{proposition}

({\v C}adek--Van{\v z}ura 1998.) The same relation holds for $\Spin(7)$-structures.

%------------------------------------------------------------
\section{Positive Quaternion-K\"{a}hler Manifolds and Wolf Spaces}
%------------------------------------------------------------

\begin{definition}
A quaternion-K\"{a}hler manifold $(X, g)$ is \textbf{positive} if it is geodesically complete and $R(g) > 0$.
\end{definition}

\begin{proposition}
A connected positive quaternion-K\"{a}hler manifold is necessarily compact and simply connected. (Salamon 1982.)
\end{proposition}

\begin{proposition}
For each dimension $4n$, there are only finitely many isometry classes of positive quaternion-K\"{a}hler manifolds. (LeBrun--Salamon 1994.)
\end{proposition}

\begin{example}[Wolf Spaces]
The \textbf{Wolf spaces} are the positive quaternion-K\"{a}hler symmetric spaces, one in each dimension $4n$. The archetypical example is the quaternionic projective plane $\HH P^2$. All Wolf spaces are positive quaternion-K\"{a}hler manifolds.
\end{example}

\begin{conjecture}[Wolf Space Conjecture]
In every dimension, the Wolf spaces are the only positive quaternion-K\"{a}hler manifolds.
\end{conjecture}

This conjecture has been verified in the cases:
\begin{itemize}
  \item $\dim = 4$: Hitchin,
  \item $\dim = 8$: Poon--Salamon (1991), LeBrun--Salamon (1994).
\end{itemize}

%------------------------------------------------------------
\section{Summary: Berger Classification for $\HH$-Geometry}
%------------------------------------------------------------

\begin{center}
\renewcommand{\arraystretch}{1.5}
\begin{tabular}{lllc}
\toprule
\textbf{Manifold} & \textbf{Holonomy} & \textbf{Dimension} & \textbf{Scalar Curvature} \\
\midrule
Hyperk\"{a}hler & $\Sp(n)$ & $4n$ & $R = 0$ \\
Quaternion-K\"{a}hler & $\Sp(n)\cdot\Sp(1)$ & $4n$ & $R = \mathrm{const}$ \\
Positive qK (Wolf) & $\Sp(n)\cdot\Sp(1)$ & $4n$ & $R > 0$ \\
\bottomrule
\end{tabular}
\end{center}

%------------------------------------------------------------
\section*{References}
%------------------------------------------------------------

\begin{itemize}
  \item Salamon, S., \textit{Quaternionic K\"{a}hler manifolds}, Invent.\ Math.\ \textbf{67} (1982), 143--171.
  \item Beauville, A., \textit{Vari\'{e}t\'{e}s K\"{a}hleriennes dont la premi\`{e}re classe de Chern est nulle}, J.\ Diff.\ Geom.\ \textbf{18} (1983), 755--782.
  \item Poon, Y.S.\ and Salamon, S., \textit{Quaternionic K\"{a}hler 8-manifolds with positive scalar curvature}, J.\ Differential Geom.\ \textbf{33} (1991), 363--378.
  \item LeBrun, C.\ and Salamon, S., \textit{Strong rigidity of positive quaternion K\"{a}hler manifolds}, Invent.\ Math.\ \textbf{118} (1994), 109--132.
  \item O'Grady, K., \textit{Desingularized moduli spaces of sheaves on a K3}, J.\ Reine Angew.\ Math.\ \textbf{512} (1999).
  \item O'Grady, K., \textit{A new six dimensional irreducible symplectic variety}, J.\ Algebraic Geom.\ \textbf{12} (2003).
  \item Joyce, D., \textit{Compact Manifolds with Special Holonomy}, Oxford University Press, 2000.
  \item Besse, A., \textit{Einstein Manifolds}, Springer, 1987.
  \item \v{C}adek, M.\ and Van\v{z}ura, J., \textit{Almost quaternionic structures on eight-manifolds}, Osaka J.\ Math.\ \textbf{35} (1998).
  \item nLab: \href{https://ncatlab.org/nlab/show/Kähler+manifold}{K\"{a}hler manifold}, \href{https://ncatlab.org/nlab/show/hyperkähler+manifold}{hyperk\"{a}hler manifold}, \href{https://ncatlab.org/nlab/show/quaternion-Kähler+manifold}{quaternion-K\"{a}hler manifold}.
\end{itemize}

\end{document}
